\section*{A three-state table, taken apart}
\addcontentsline{toc}{section}{A three-state table, taken apart}

The decomposition becomes concrete in the smallest example that retains more than one categorical
contrast.  Let both endpoints have three states and use uniform references
\begin{equation}
p_i=p_j=\frac13(1,1,1)^{\mathsf T}.
\end{equation}
Consider the raw pair table
\begin{equation}
T=
\begin{pmatrix}
10&5&3\\
5&5&2\\
6&5&4
\end{pmatrix}.
\label{eq:book-three-state-raw}
\end{equation}
Its grand mean is $c=5$.  Its row means are $(6,4,5)$, so the centered target-only term is
\begin{equation}
u=(1,-1,0)^{\mathsf T}.
\end{equation}
Its column means are $(7,5,3)$, giving the centered source-only term
\begin{equation}
v=(2,0,-2)^{\mathsf T}.
\end{equation}
Subtracting these three lower-order pieces leaves
\begin{equation}
G=T-c\mathbf1\mathbf1^{\mathsf T}
-u\mathbf1^{\mathsf T}-\mathbf1v^{\mathsf T}
=
\begin{pmatrix}
2&-1&-1\\
-1&1&0\\
-1&0&1
\end{pmatrix}.
\label{eq:book-three-state-gauged}
\end{equation}
Every row and column of $G$ sums to zero.  Equivalently,
\begin{equation}
p_i^{\mathsf T}G=0,
\qquad
Gp_j=0.
\end{equation}
The large entry $T_{11}=10$ is therefore not itself the interaction.  It contains the grand
baseline, a target propensity, a source propensity, and the bivariate residual.  In this example
those contributions are $5+1+2+2$.

Now alter the raw table by adding any new row and column effects,
\begin{equation}
T'(a,b)=T(a,b)+\widetilde u(a)+\widetilde v(b)+\widetilde c.
\end{equation}
The numerical entries and their ordinary Frobenius norm can change substantially, but double
projection returns the same $G$.  The identified interaction is therefore the quotient class of
$T$, while Eq.~\ref{eq:book-three-state-gauged} is its representative in the uniform reference.

The edge becomes an operator when it acts on a source contrast.  Take
\begin{equation}
f=(1,-1,0)^{\mathsf T},
\qquad p_j^{\mathsf T}f=0.
\end{equation}
Then
\begin{equation}
\mathcal K_{i\leftarrow j}f
=GD_jf
=\frac13Gf
=
\begin{pmatrix}
1\\[-2pt]
-2/3\\[-2pt]
-1/3
\end{pmatrix},
\label{eq:book-three-state-action}
\end{equation}
whose target mean is again zero.  The source contrast does not merely activate an edge of strength
$\|G\|$; it produces a particular target contrast.  Another source direction produces another
response.

For uniform references the whitened edge matrix is $\widetilde G=G/3$.  Its two nonzero singular
values are $1$ and $1/3$, and its Hilbert--Schmidt norm is
\begin{equation}
\|\mathcal K_{i\leftarrow j}\|_{\mathrm{HS}}
=\frac{\sqrt{10}}{3}.
\end{equation}
The scalar norm retains their combined magnitude but discards both singular directions.  This one
calculation contains the sequence of transformations used throughout the book:
\begin{equation}
\boxed{
\text{raw table}
\longrightarrow
\text{gauge class}
\longrightarrow
\text{reference representative}
\longrightarrow
\text{fiber operator}
\longrightarrow
\text{scalar shadow}.}
\label{eq:book-table-to-shadow}
\end{equation}
Only the first four objects retain the typed substitution action.  The final scalar is useful only
after the loss has been declared.
